\section{Basic requirements}
\label{sec:prb_requirements}
The challenges outlined in the previous section are not exhaustive, nor are they entirely independent of each other. They are, however, representative of the issues that typically arise when designing an autonomous system. Summarizing briefly the many aspects previously covered, suppose that the task at hand is that of developing an autonomous system capable of performing a series of operations. It is reasonable, in the context of this work, to focus on situations that meet the following assumptions.
\begin{itemize}
  \item The system has a physical structure that can be modeled as a dynamic system, thus, one of the tasks consists in stabilizing and controlling this dynamic system, something that is commonly referred to in the following as \emph{low-level control} or simply \emph{control}.
  \item Digital programmable computers of different kinds are used to both develop and run all the algorithms, ranging from low-level control to task planning and supervision.
\end{itemize}
Note that the situations hereby described might span from mobile robots to drones and even power plants, all classes of systems that have been subject of this research work. Recalling \Cref{sec:prb_design_challenges}, a series of choices regarding the following items have to be made:
\begin{itemize}
  \item sensors and actuators to use;
  \item software tools and packages to employ and integrate;
  \item hardware platforms, \emph{i.e.}, onboard electronics to adopt.
\end{itemize}
Note how these items might be strongly interconnected. Low-level hardware has to be chosen according to the control objective for the underlying dynamic system, while high-level hardware is selected usually weighing the software that it runs against the physical capabilities, as well as the costs, of the system that is going to host it. Finally, software tools and packages can be divided in two groups: that which has to be developed from scratch, possibly very specific to the system at hand, and that which can be reused from existing libraries or frameworks. The latter category, for example, includes solutions for data processing, communication, and visualization, as well as remote access and monitoring tools; when one wants to develop the prototype of an autonomous system, especially for research purposes, these things may be overlooked but they are crucial for the system to be usable and maintainable and the overhead that they might introduce can easily become significant if not properly managed early on.

There is a principle in computer science \cite{raymond_how_2001}, dating back to the very early days of the discipline, stating that to do a proper job \emph{no problem should ever be solved twice} or, in a historical fashion, \emph{one must never reinvent the wheel} with the only exception of building a better wheel. One can see how this fits perfectly in the context of the development of autonomous systems: there is a wide range of complex foundational and integration problems to solve and present themselves each time a new system is developed. The scope of this work is to provide a solution that is \emph{reusable}, \emph{modular}, and \emph{maintainable} and that can be used as a foundation for the development of innovative prototypes and products. With this in mind, the following requirements are formally identified.
\begin{description}
  \item[(M) Modularity] The architecture should constitute both a development and deployment framework for solutions ranging from single modules, meant to be easily integrated and reused, to entire autonomous systems prototypes. To guarantee this, the structure of the framework should suit both these use cases: single modules should carry with them all their dependencies, and it should be possible to easily \emph{merge} single modules to build a \emph{superproject} out of many \emph{subprojects}.
  \item[(D) Development] The architecture should support the development of single modules both as standalone packages and as parts of the software suite of a prototype. While working on a superproject that integrates a module, it should be possible to modify and update the module easily propagating the updates to every other superproject which integrates that module, unless the superproject has been specifically configured to use a different version of the module. The architecture should also ease testing and debugging activities, independently of the environment in which a module is developed.
  \item[(A) Abstraction] The architecture should offer a \emph{system abstraction layer}, \emph{i.e.}, it should make it possible to carry on all the development activities without having to care too much about the underlying hardware on which a module or prototype software suite is deployed, apart from unavoidable specific differences, as well as to the target OS of such a platform. The architecture should ease deployment by offering a common framework across multiple different hardware platforms. In the end, the architecture shall offer very similar, if not equal, development and execution environments across different hardware platforms, requiring minimal effort to be replicated.
  \item[(O) Overhead] The architecture should be as lightweight as possible, in terms of both computational and storage resources. This is to ensure that the overhead introduced by the framework itself is minimal, and that the resources of the hardware platforms on which the framework is deployed can be used to their full potential for the tasks that the system is meant to perform.
\end{description}
Note that these requirements, if examined by the perspective of subjects such as systems engineering, might not appear well-posed due to, \emph{e.g.}, the lack of a clear and objective way to evaluate them. They are, instead, meant to guide the design of an architecture and not that of a specific system. Moreover, the design stage at which they are introduced is one in which many things still have to be evaluated such as, \emph{e.g.}, hardware platforms to support and preexisting software tools to integrate, and thus the final definition of the architecture is the result of many trade-offs. They are high-level requirements meant to guide the design process, characterizing the solution from a broad and functional point of view.
